% ---------------------------------------------------------------------
%  HERRAMIENTAS MATEMÁTICAS Y ENTORNOS
% ---------------------------------------------------------------------
\usepackage{bm}
\usepackage{amsthm}

% --- estilo "plain": encabezado en negrita, cuerpo en cursiva ---
\newtheoremstyle{plainbook}
  {12pt plus 2pt minus 2pt}   % espacio antes
  {12pt plus 2pt minus 2pt}   % espacio después
  {\itshape}                  % cuerpo
  {}                          % sangría
  {\bfseries}                 % encabezado
  {.}                         % puntuación tras el encabezado
  {0.6em}                     % espacio tras el encabezado
  {}

% --- estilo "definición": encabezado en negrita, cuerpo en romana ---
\newtheoremstyle{defbook}
  {12pt plus 2pt minus 2pt}
  {12pt plus 2pt minus 2pt}
  {\normalfont}
  {}
  {\bfseries}
  {.}
  {0.6em}
  {}

% --- estilo "observación/nota": encabezado en cursiva, cuerpo romano ---
\newtheoremstyle{notebook}
  {10pt plus 2pt minus 2pt}
  {10pt plus 2pt minus 2pt}
  {\normalfont}
  {}
  {\itshape}
  {.}
  {0.6em}
  {}

% Cada entorno recibe su propio contador (para que cleveref use el nombre
% correcto de cada tipo), pero todos se alían al contador de "theorem".
\theoremstyle{plainbook}
\newtheorem{theorem}{Teorema}[chapter]

\theoremstyle{plainbook}
\newtheorem{lemma}{Lema}[chapter]
\newtheorem{proposition}{Proposición}[chapter]
\newtheorem{corollary}{Corolario}[chapter]
\newtheorem{conjecture}{Conjetura}[chapter]

\theoremstyle{defbook}
\newtheorem{axiom}{Axioma}[chapter]
\newtheorem{definition}{Definición}[chapter]
\newtheorem{example}{Ejemplo}[chapter]
\newtheorem{exercise}{Ejercicio}[chapter]
\newtheorem{problem}{Problema}[chapter]

\theoremstyle{notebook}
\newtheorem{remark}{Observación}[chapter]
\newtheorem{notation}{Notación}[chapter]

\makeatletter
\let\c@lemma\c@theorem
\let\c@proposition\c@theorem
\let\c@corollary\c@theorem
\let\c@conjecture\c@theorem
\let\c@axiom\c@theorem
\let\c@definition\c@theorem
\let\c@example\c@theorem
\let\c@exercise\c@theorem
\let\c@problem\c@theorem
\let\c@remark\c@theorem
\let\c@notation\c@theorem
\makeatother

% Demostración con marca cuadrada discreta
\renewcommand{\qedsymbol}{$\mdlgblksquare$}
\renewcommand{\proofname}{Demostración}
